We did the second assignment presentation later this week and got feedback from the instructors. To me the main point from the instructors is that our product is less "enchanted" or magical than our previous snow globe draft. We started to revise our design at the beginning of the low fidelity prototype phrase. We shifted our focus back to a snow globe rather than a lamp, as the feedback suggesting that a snow globe could also be considered an everyday object. During the discussion, I referred to the user research results that I mentioned earlier, including these main points: 
\begin{enumerate}
    \item We could narrow down our use cases. Currently we plan to cover family, friendship, and couple relationship, but participants point out some considerable differences and unique characteristics.
    \item The voice message function should be either complete (e.g. phone call) or just be replaced by another function. The experience of our currently asynchronous voice message function would not be good enough. They also give some examples.
    \item The tap function is not rich in social cues, and (same as 2) asynchronous voice messages are not ideal for the next move after they reply to others’ tap immediately
    \item Participants also brainstorm another similar idea of digital album with colourful e-ink display.
\end{enumerate}

Although we did not applied every of them and still continue our snow globe plan, we all agreed to narrow the use case to one-on-one interactions. We also simplified the feature set to include only voice messaging, shaking the snow globe to convey a signal, and displaying ambient lighting and weather information. The enchantment element primarily arises from the turning to play voice messages and the shaking to send a magical signal function. We envisioned allowing the user to turn the snow globe to receive voice messages and shake it to activate a miniature storm inside, effectively combining the snow globe with a music box.

 The build of the low fidelity prototype happened simultaneously when we revise our design. We initially built a turnable platform using LEGO blocks to demonstrate the turning gesture. Later, we decided to use the provide sticky foam to improve our prototype, but this proved time-consuming without adding more useful information. Therefore, we reverted to the LEGO prototype. Personally, I was also studying 3D modeling and printing at the same time. In this stage, I printed a clockwork mechanism to provide mechanical feedback for the turning operation. For the low-fi prototype, we still needed a transparent shell. I experimented with transparent PETG material for 3D printing, attempting to print a single-layer shell in vase mode, but it failed because the top area lacked proper support. Choosing non-vase mode resulted in an opaque outcome. Despite multiple experiments, I eventually gave up. Additionally, Marit and Otto kindly showed us the glass lab at Utrecht University, which was truly impressive. Anticipating a glass shell for the high-fi prototype, we stopped trying to print a transparent shell for the low-fi prototype. Instead, we used an opaque one for users to experience the turning operation and started working on the high-fi prototype.

As for the design iteration, we elaborate our design choices to enchant a snow globe. I will mention it in the later high fidelity prototype section.